\frontmatter
\begin{titlepage}
\centering
\vspace*{1.4cm}
{\sffamily\bfseries\fontsize{36}{40}\selectfont\color{croqink} Croqtile\par}
\vspace{0.1cm}
{\sffamily\bfseries\fontsize{26}{30}\selectfont\color{croqgreen} 编程指南\par}
\vspace{0.8cm}
{\Large\itshape 面向高性能核函数的 Croqtile 语言导读\par}
\vspace{1.1cm}
\rule{0.72\linewidth}{0.6pt}\par
\vspace{0.9cm}
{\large 版本 \guideversion{} --- \guidereleasedate\par}
\vfill
{\small\color{croqgray} LancerLab\par}
\vspace{1.0cm}
\end{titlepage}
\thispagestyle{empty}
\cleardoublepage

\chapter*{前言}

\addcontentsline{toc}{chapter}{前言}

大模型的广泛应用，让算子的高效实现成为一项越来越重要的开发工作。模型架构逐渐形成了几类常见设计，内部的计算却仍在不断变化：注意力机制、稀疏模式、位置编码和量化方案的演进，都带来了新的实现需求。即使是同一个算子，面对不同的输入形状和数值格式，也可能需要不同的核函数(kernel)。模型采用相近的架构，底层仍有大量实现和优化工作需要完成。

硬件的多样化进一步增加了这项工作的规模。同一套程序可能需要运行在多种加速器上，而它们的存储组织和执行方式各不相同；即便属于同一系列，新一代硬件引入的指令和编程接口，也可能要求开发者重新审视并调优已有的核函数。算子融合又带来了更多实现组合：把多个操作合并，可以减少中间数据搬运(data movement)和启动开销，但每种组合都有自己的实现选择。要支持这些计算需求，就需要持续开发新的核函数，也需要更高效的开发方式。

写好这些核函数并不容易。分块(tile)尺寸、数据布局、数据暂存方式、流水线深度以及并行区域之间的协作，共同决定了性能；其中一项发生变化，往往还要调整其他几项。近年来陆续出现的高层核函数语言，正是希望用更少的代码表达这些工作，提高开发效率。它们面临的核心问题是：如何简化编程，同时让开发者仍能做出实现高性能所需的决策。如果抽象隐藏了这些决策，能够探索的实现方案也会随之受限。

Croqtile 是一门面向高性能异构计算的编程语言，设计的出发点是提高开发者的效率，让核函数更容易编写和修改。形状、分块选择、数据搬运与协调都有直接的语法表达，开发者可以用一段紧凑的程序描述核函数的整体工作方式。相关的尺寸和索引由带形状的数据推导，减少了修改程序结构时需要同步维护的细节；符号维度(symbolic dimension)让同一份源码适用于不同的输入尺寸，编译时检查则能在运行前发现形状和索引的不一致。这些抽象不增加运行时开销，也让影响性能的决策清楚地保留在程序中。

同样的设计也为 AI 辅助编程与调优提供了支持。源码更紧凑，智能体的上下文就能容纳更多性能分析结果和历史尝试；各项操作可以独立表达，调整流水线或并行组织时，只需修改少数几处代码；编译器的诊断则帮助它定位并修正不合法的组合。智能体由此能够反复提出修改、检查程序、测量效果，再据此继续改进。面向 AI 的设计，把提高核函数开发效率的目标延伸到了自动化工作流中。

Croqtile 嵌入在 C++ 中。在 \co@.co@ 文件里，普通 C++ 负责准备带形状的缓冲区并调用 \co@__co__@ 函数，后者描述并行工作的迭代、分块、数据搬运与协调。本指南通过示例介绍这些构造，并说明如何将它们组合成完整程序。

\chapter*{阅读指南}

\addcontentsline{toc}{chapter}{阅读指南}

本指南面向希望读懂并编写 Croqtile 核函数的开发者。阅读前，只需掌握 C++ 中的函数、指针与数组，并了解张量的形状。书中的示例从基本构造讲起，逐步说明如何组织并行计算、选择与搬运分块，以及协调各部分的执行。你可以跟随示例读懂 \co@.co@ 程序，再把这些构造组合起来，编写自己的核函数。

如果你第一次接触 Croqtile，建议先阅读第 1 章“入门”，完成环境准备、编译和运行，确认示例结果正确，再进入第 2 章“语言基础”。它通过简短的程序，从函数声明和形状描述开始，逐步加入迭代、分块、数据搬运与协调。阅读时，可以先看程序的整体结构，再对照说明理解每条语句的作用，重点把握数据怎样划分、搬运和参与计算。

熟悉基本写法后，可以按需要查阅第 3 章“编程指南”。它按功能组织内容，进一步解释各类构造的语义、使用约束与设计理由。完整核函数与相关特性放在一起，既便于深入理解某个主题，也方便在编写程序时随时查阅。附录 A 集中列出语法，适合快速确认写法；需要理解某个构造的含义和适用条件时，再回到相应章节。
